\subsection{Cross-Dataset and Data-Regime Comparisons}
\label{sec:extended_benchmarks}

The controlled CIFAR experiments isolate allocation at a fixed dropout budget; the broader benchmark asks how the resulting training recipes behave across finance, tabular classification, text, audio, and vision. We compare depth-12, width-256 MLPs and Transformers, tune learning rate and mean dropout separately for each profile using validation data, and evaluate fresh confirmation seeds at their minimum-validation-loss checkpoints. Table~\ref{tab:benchmark_results} reports the uniform baseline and best observed nonuniform profile for every recovered cohort, with complete per-profile endpoints and learning curves in App.~\ref{app:extended_experiments}. Because the mean probability is tuned separately, these comparisons measure the resulting recipes rather than holding allocation as the only varying factor; the winner is selected descriptively from the observed test means, so the reported improvements do not constitute an independent test of a prespecified schedule.

In the zero-weight-decay cohorts, the best nonuniform profile lowers mean test cross-entropy in all ten dataset--architecture cells, with reductions from approximately $0.5\%$ for the Jannis Transformer to $12.7\%$ for the Speech Commands MLP. The ordering is architecture dependent: early profiles attain the lowest observed loss in the five MLP cells, whereas increasing linear profiles attain it for the FI-2010 and Amazon Transformers. The data-regime comparisons further delimit the effect: uniform dropout has lower test loss than every completed nonuniform candidate for the Amazon Transformer at both 2,000 and 20,000 training examples, and the 100-epoch Jannis follow-up is effectively tied at the reported precision. These outcomes motivate treating spatial allocation as a task- and architecture-dependent regularization choice, while the original fixed-budget experiments provide the more direct tests of the mean-field scheduling mechanism.
